\section{Related Work}
\label{sec:related}

\paragraph{Classical models for tabular data.}
Logistic regression \citep{cox1958logistic} and least squares remain strong baselines for small, low-dimensional
problems. Random forests \citep{breiman2001rf} and extremely randomized trees \citep{geurts2006extremely} average
decorrelated trees; gradient boosting \citep{friedman2001gbm} fits trees sequentially to residuals, and its efficient
implementations XGBoost \citep{chen2016xgboost}, LightGBM \citep{ke2017lightgbm} and CatBoost
\citep{prokhorenkova2018catboost} are the usual reference standard. We use the scikit-learn implementations
\citep{pedregosa2011sklearn} where available.

\paragraph{Deep learning on tables.}
Comparisons of neural architectures such as TabNet \citep{arik2021tabnet} and tabular transformers
\citep{gorishniy2021revisiting} with tree ensembles found trees competitive or superior on many medium-sized datasets
\citep{shwartzziv2022tabular,borisov2024survey,grinsztajn2022tree}. Well-regularized multilayer perceptrons can close
the gap \citep{kadra2021welltuned}, and the best family varies with dataset properties \citep{mcelfresh2023neural}.

\paragraph{In-context learning and tabular foundation models.}
In-context learning, in which a pretrained model solves a task from examples in its input without weight updates,
was popularized by large language models \citep{brown2020gpt3} built on transformers \citep{vaswani2017attention},
and emerges for simple function classes under suitable training distributions \citep{garg2022icl}. Language models
have also been applied to serialized tables \citep{hegselmann2023tabllm}. Prior-data fitted networks
\citep{muller2022pfn} are trained on datasets sampled from a prior so that a forward pass approximates the posterior
predictive distribution. TabPFN \citep{hollmann2023tabpfn} applied this to small classification tables, and its
successor \citep{hollmann2025tabpfn} extended it to regression and larger inputs. TabICL \citep{qu2025tabicl}
compresses rows with column- and row-wise attention before in-context learning; set attention with inducing points
\citep{lee2019set} is a common component.

\paragraph{Kumo Tabular.}
Kumo Tabular \citep{nvidia_kumo_card,nvidia_kumo_blog} follows this line. Its developers describe column, row and
in-context attention ``as introduced in TabICL and TabPFN'', pretraining on synthetic tables sampled from structural
causal models, and three sizes with 28M to 215M parameters; they report leading results on several benchmark suites.
Those claims are the developers' own. We are not aware of an independent evaluation of Kumo Tabular.

\paragraph{Benchmarking methodology.}
OpenML \citep{vanschoren2014openml} provides datasets with fixed evaluation splits; OpenML-CC18
\citep{bischl2021openml} and OpenML-CTR23 \citep{fischer2023ctr23} are curated classification and regression suites.
For comparisons across datasets we follow \citet{demsar2006comparisons}: a Friedman test \citep{friedman1937ranks}
over all methods, with pairwise Wilcoxon signed-rank tests \citep{wilcoxon1945} and Holm's correction
\citep{holm1979}. For comparisons within one dataset over resampled splits, \citet{dietterich1998tests} showed that
naive paired tests are anti-conservative, and \citet{nadeau2003inference} proposed the corrected resampled $t$-test
we use for the secondary within-dataset analysis.
